\documentclass[10pt,a4paper]{amsart}
\usepackage[margin=19mm]{geometry}
\usepackage{amsmath,amssymb,mathtools}
\usepackage[scr=rsfs,cal=euler]{mathalfa}
\usepackage{xfrac}
\usepackage[unicode,hidelinks,bookmarks=false]{hyperref}
\newcommand{\ie}{i.e.\ }
\newcommand{\cf}{cf.\ }
\newcommand{\resp}{respectively\ }
\newcommand{\Subsection}{\subsection}
\providecommand{\nobreakdash}{\nobreak\hbox{-}}
\setlength{\emergencystretch}{2em}
\multlinegap0pt
\usepackage{mathtools}
\usepackage{tensor}
\usepackage[normalem]{ulem}
\usepackage[shortlabels]{enumitem}
\newtheorem{prop}{Proposition}
\newcommand{\Fb}{\mathbb{F}}
\newcommand{\Pb}{\mathbb{P}}
\title{Exact counts of elliptic curves of bounded height over $\mathbb F_q(t)$ in characteristics $2$ and $3$}
\author{Jun-Yong Park}
\date{Source: arXiv:2507.06754v3 (CC BY 4.0)}
\begin{document}
\maketitle

\noindent\emph{Source and licence.} Exact counts of elliptic curves of bounded height over $\mathbb F_q(t)$ in characteristics $2$ and $3$. Version arXiv:2507.06754v3 (CC BY 4.0), distributed by arXiv under the Creative Commons Attribution 4.0 International licence, http://creativecommons.org/licenses/by/4.0/, as recorded in the arXiv licence metadata for this exact version and shown on the abstract page https://arxiv.org/abs/2507.06754v3. Full licence notice: \texttt{licenses/arxiv-2507.06754-CC-BY-4.0.txt}.\par\smallskip

\noindent\textit{Editorial scope.} Counted unit: the article's prop prop:nonsmooth (source0.tex lines 642--651). The article's complete original proof of it is delivered with it (source0.tex lines 653--752, one \texttt{proof} environment of 100 lines). No mathematics is written, repaired or re-derived: the statement and the proof are the article's own lines, in the article's own notation, and no retained line is substituted or re-flowed.  No further source-local context is needed: every cross-reference of every retained line resolves inside the two delivered spans.  Nothing was omitted from any retained span, so the package carries no omission record.  No retained line contains a citation, so the wrapper carries no bibliography and no bibliography entry.  No retained line contains a figure environment, an image-inclusion command or drawing code, so no image is delivered and the package declares no asset dependency.  Editorial formatting only, in addition: an AMS \texttt{amsart} preamble that restates the article's own declarations and macros used by the retained lines, namely the environments \texttt{prop} on separate counters, together with the standard packages those lines need.  The retained environments print in this document as Proposition 1 in order of appearance, whereas the article prints the same items with the numbering of its own full document.  The complete native source of the article as distributed in the arXiv e-print for this exact version is bundled unchanged as \texttt{sources/181286/source0.tex}.
\begin{prop}\label{prop:nonsmooth}
Let $p$ be a prime, let $r\geq1$, put $q=p^r$, and let $d\geq0$.  If $p\in\{2,3\}$, the number of coefficient-degree-$d$ generalized Weierstrass equations over $\Pb^1_{\Fb_q}$ whose generic fiber is not smooth equals
\[
\begin{cases}
q^{8d+4}+q^{10d+3}-q^{6d+3},&p=3,\\[1mm]
q^{8d+4}+q^{12d+3}-q^{7d+3},&p=2.
\end{cases}
\]
If $p\geq5$, the number is $q^{8d+4}$.  In particular, for $d=0$ the answer is $q^4$ in every characteristic.
\end{prop}

\begin{proof}
Write the equation as $F=0$, where
\[
F
=
y^2+a_1xy+a_3y-x^3-a_2x^2-a_4x-a_6,
\qquad
a_i\in H^0\bigl(\Pb^1,\mathcal O(id)\bigr).
\]
The generic fiber is smooth if and only if its discriminant is nonzero in $K$.

We first count the equations whose generic fiber has a $K$-rational singular point.  After homogenization, the generic Weierstrass cubic meets the line at infinity only at $O=[0:1:0]$.  This point is smooth, since the derivative of the homogenized equation with respect to the homogenizing variable equals $1$ at $O$.  Every positive-degree projective component meets the line at infinity, so every geometric component of the cubic passes through $O$.  If the cubic were reducible, at least two components would pass through $O$; if it were nonreduced, a repeated component would pass through $O$.  Either possibility would make $O$ singular.  Thus the cubic is geometrically integral.  By B\'ezout's theorem, a geometrically integral plane cubic has at most one geometric singular point: otherwise, the line joining two distinct singular points would meet the cubic with total intersection multiplicity at least four.

By the valuative criterion of properness over the regular curve $\Pb^1$, a $K$-rational singular point extends uniquely to a section of the associated Weierstrass surface.  The relative singular locus, namely the closed locus where the morphism to $\Pb^1$ is not smooth, contains the generic point of this section and hence the entire section.  The zero section lies in the relative smooth locus, and every fiber meets the line at infinity only in the zero section.  Hence the singular section lies in the affine chart and has the form
\[
(x,y)=(C_2,C_3),
\qquad
C_2\in H^0\bigl(\Pb^1,\mathcal O(2d)\bigr),
\qquad
C_3\in H^0\bigl(\Pb^1,\mathcal O(3d)\bigr).
\]
At $(C_2,C_3)$, the equations $F=F_x=F_y=0$ are equivalent to
\[
\begin{aligned}
a_3&=-2C_3-a_1C_2,\\
a_4&=a_1C_3-3C_2^2-2a_2C_2,\\
a_6&=-C_3^2-a_1C_2C_3+2C_2^3+a_2C_2^2.
\end{aligned}
\]
Since the geometric singular point is unique, $(a_1,a_2,C_2,C_3)$ parametrizes these equations bijectively.  Their number is
\[
q^{(d+1)+(2d+1)+(2d+1)+(3d+1)}=q^{8d+4}.
\]

We use the following elementary observation.  Let $m\geq1$ and let $s$ be a rational section of a line bundle $L$ on $\Pb^1$.  If $s^m\in H^0(\Pb^1,L^{\otimes m})$, then $s$ is global: in a local trivialization at every closed point $x$, one has $m v_x(s)=v_x(s^m)\geq0$.

Suppose first that $p=2$.  On the affine chart, $\partial F/\partial y=a_1x+a_3$.  If $a_1\neq0$, every geometric singular point satisfies
\[
x_0=\frac{a_3}{a_1},
\qquad
a_1y_0=x_0^2+a_4,
\]
and is therefore $K$-rational.  If $a_1=0$ and $a_3\neq0$, then $\partial F/\partial y=a_3$ is nonzero on the generic fiber, so the generic fiber is smooth.  If $a_1=a_3=0$, the singularity equations over $\overline K$ are
\[
x_0^2=a_4,
\qquad
y_0^2=a_2a_4+a_6.
\]
They have a unique solution in $\overline K^2$, so all $q^{(2d+1)+(4d+1)+(6d+1)}=q^{12d+3}$ equations in this stratum are nonsmooth.

Within this stratum, the singular point is $K$-rational precisely when there are rational sections $C_2$ of $\mathcal O(2d)$ and $C_3$ of $\mathcal O(3d)$ such that
\[
a_4=C_2^2,
\qquad
a_6=C_3^2+a_2C_2^2.
\]
The preceding observation shows that both sections are global.  Since the squaring maps on their spaces of sections are injective, the triples $(a_2,C_2,C_3)$ parametrize exactly $q^{(2d+1)+(2d+1)+(3d+1)}=q^{7d+3}$ members of this stratum already counted by the rational-singular parametrization.  Inclusion--exclusion therefore gives
\[
q^{8d+4}+q^{12d+3}-q^{7d+3}.
\]

Now suppose that $p=3$.  Completing the square removes $a_1$ and $a_3$. For each fixed pair $(a_1,a_3)$, this is an affine-linear bijection on the remaining coefficient space, and there are $q^{(d+1)+(3d+1)}=q^{4d+2}$ choices of $(a_1,a_3)$.  Renaming the transformed coefficients gives
\[
y^2=f(x),
\qquad
f=x^3+a_2x^2+a_4x+a_6.
\]
A singular point satisfies
\[
y_0=0,
\qquad
f(x_0)=f'(x_0)=0,
\qquad
f'=2a_2x+a_4.
\]
If $a_2\neq0$, every repeated root is $x_0=-a_4/(2a_2)\in K$.  If $a_2=0$ and $a_4\neq0$, then $f'=a_4$ is nonzero, so the generic fiber is smooth.  If $a_2=a_4=0$, then $f=x^3+a_6$.  For each fixed pair $(a_1,a_3)$, all $q^{6d+1}$ choices of $a_6$ in this stratum are nonsmooth.  The singular point is $K$-rational precisely when $a_6=-C_2^3$ for a rational section $C_2$ of $\mathcal O(2d)$.  The preceding observation shows that $C_2$ is global, and injectivity of the cubing map gives $q^{2d+1}$ such equations.

For a rational singular point with $x=C_2$, the equations $f'(C_2)=f(C_2)=0$ give
\[
a_4=-2a_2C_2,
\qquad
a_6=-C_2^3+a_2C_2^2.
\]
For each fixed $(a_1,a_3)$, the rational-singular locus has $q^{4d+2}$ elements, the stratum $a_2=a_4=0$ has $q^{6d+1}$ elements, and their intersection, obtained by setting $a_2=0$, has $q^{2d+1}$ elements. Restoring the $q^{4d+2}$ choices of $(a_1,a_3)$, inclusion--exclusion gives
\[
q^{4d+2}
\bigl(q^{4d+2}+q^{6d+1}-q^{2d+1}\bigr)
=q^{8d+4}+q^{10d+3}-q^{6d+3}.
\]

For $p\geq5$, completing the square and depressing the cubic gives
\[
y^2=f(x),
\qquad
f=x^3+Ax+B.
\]
These coordinate changes are defined over $K$ and preserve the rationality of singular points.  A singular point has $y=0$ and is a common root of $f$ and $f'=3x^2+A$.  If $A=0$, the only root of $f'$ is $0\in K$.  If $A\neq0$, then $f'$ is separable.  An irrational common root of $f$ and $f'$ would have a distinct conjugate that was also a common root, giving two distinct repeated roots of a cubic, which is impossible. Hence every nonsmooth equation has a $K$-rational singular point, and its number is $q^{8d+4}$.

Finally, when $d=0$, Frobenius is bijective on $\Fb_q$, so the small-characteristic strata add no equations beyond the rational-singular locus.  Both formulas reduce to $q^4$.
\end{proof}

\end{document}
